% =====================================================================
% 《AI 工具使用详情》独立编译源文件
% 编译方式：xelatex AI工具使用详情.tex
% 生成 PDF 后请按官方要求将文件名改为：AI工具使用详情.pdf
% AI 使用详情单独提交，不显示参赛队号。
% =====================================================================
\documentclass[UTF8,a4paper,12pt]{ctexart}
\usepackage[left=2.5cm,right=2.5cm,top=2.5cm,bottom=2.5cm]{geometry}
\usepackage{amsmath,amssymb,bm}
\usepackage{booktabs,tabularx,array}
\usepackage{listings}
\usepackage[colorlinks,linkcolor=blue]{hyperref}

% 代码块风格：简单四方框。columns=fullflexible + keepspaces 仅用于
% 保证中文注释与缩进正常显示，不属于任何视觉装饰。
\lstdefinestyle{ai}{%
  basicstyle=\ttfamily\small,
  frame=single,
  breaklines=true,
  columns=fullflexible,
  keepspaces=true,
  showstringspaces=false,
}

\title{\textbf{AI 工具使用详情}}
\author{}
\date{2026 年 9 月}

\begin{document}

\maketitle
\vspace{-1em}

%--------------------- 0. 引言 ---------------------
\section{撰写说明}

本题为 2026 年高教社杯全国大学生数学建模竞赛 \textbf{B 题：无线电干扰源
的快速自动定位与清除}。本团队严格按照竞赛关于生成式人工智能的使用规范
要求，如实记录使用 AI 辅助完成本赛题的过程，内容包括：所用 AI 工具名称
与型号、具体使用目的与环节、主要提示方式与使用过程说明（附典型交互
示例），以及对 AI 输出的采纳、人工修改与核验情况，语言润色除外。

需要首先声明 AI 的作用边界：\textbf{本论文所有核心模型的构思、数学
推导、算法设计与整体技术方案均由团队独立完成}。AI 仅被用于算法实现
细节、公式/代码的排错与校核、代码调试与可视化脚本撰写、以及论文结构
梳理等辅助性环节。本说明与论文内容一致，且过程可复现、可复核。

%--------------------- 1. 工具及型号 ---------------------
\section{所用 AI 工具的名称、版本/型号}

\subsection{DeepSeek（DeepSeek-V4-flash 与 DeepSeek-V4-pro）}
\begin{itemize}
  \item 访问方式：官网网页版，网址 \texttt{chat.deepseek.com}。
  \item 型号分工：代码实现、快速问答与日常排错使用 DeepSeek-V4-flash，
        涉及公式多步推导、边界情形论证时使用推理能力更强的
        DeepSeek-V4-pro。
  \item 用途侧重：代码生成与调试、算法数值细节、边界情况梳理，
        以及部分 LaTeX 数学公式书写与核验。
\end{itemize}

\subsection{ChatGPT（OpenAI，GPT-5.6 SOL）}
\begin{itemize}
  \item 使用方式：OpenAI 网页版，主用模型 GPT-5.6 SOL，支持代码与文字
        混合任务及多轮检索。
  \item 用途侧重：论文结构与章节逻辑梳理、matplotlib 可视化脚本初稿、
        面向评审的表述改进，并与 DeepSeek 形成交叉核验。
\end{itemize}

团队全程仅把 AI 视为辅助工具，AI 生成的核心公式、算法与结论性文字均
不会直接进入论文；其输出需通过本队人工审查、改写后再决定是否采纳。

%--------------------- 2. 使用目的和环节 ----------------------
\section{具体使用目的和环节}

本赛题的四问沿“观测 → 更新可能位置 → 选择下一行动 → 清除”的行动链
展开，AI 的使用贯穿以下四个环节：

\begin{enumerate}
  \item \textbf{算法实现与代码调试}：包括问题一旋转卡壳求凸多边形最远
        点对与直径圆覆盖检验；问题二可行域投影 $d_-$、$d_+$ 与双侧候
        选点生成；问题三/四的集合成员定位、最小覆盖圆计算、最近邻开放
        路径加二边交换模拟退火，以及终止准则判断。这些是工程性、可枚举
        的实现部分，AI 帮助写出正确且较高效的 Python 实现并完成初步
        排错。
  \item \textbf{公式推导与一致性核验}：将论文中的关键公式，如示向角扇
        形对应的两个半平面叉积约束、正负观测的中垂线半平面、定向源的多
        点无信号证书区域 $\mathcal{D}_{ijk}$ 等，提交给 AI，请其作符号展开、
        检查括号与索引，核对符号 $\varepsilon'=1.0051^\circ$ 的舍入处理
        是否自洽，并给出等价再解释。
  \item \textbf{可视化与评估脚本}：撰写问题一几何示意图、问题三/四的
        轨迹重建、以及统计结果的 matplotlib 脚本初稿，随后由人工调整
        坐标、标签与配色。
  \item \textbf{论文结构梳理}：用于整理各章之间的逻辑衔接、摘要的变量
        关系措辞，以及“本文工作”的叙述结构；涉及论文最终内容与结论的
        任何措辞均由团队人工定稿。
\end{enumerate}

%--------------------- 3. 使用过程说明 ----------------------
\section{主要提示方式与使用过程说明}

\subsection{总体提示策略}
团队把复杂任务拆解后逐项提出给 AI，采用“背景 + 目标 + 边界条件与
期望的输入输出 + 约束与禁忌”的方式，并在关键点要求 AI 先给出分析思路
再给代码，以便判断其是否真正理解问题，不直接照搬输出。对可能影响结论
的关键环节，团队在决定采纳前会由核心成员人工再推导一遍。

下面给出三个典型的交互示例，保留与本论文技术内容相关的部分。

\subsubsection{示例 A：问题一 旋转卡壳求直径与直径圆覆盖判定}

\paragraph{用户提问}
\begin{lstlisting}[style=ai]
我们已知定位多边形 vertices（逆时针、凸）。需要：
1. 用旋转卡壳求最远顶点对，返回 max_dist 与 (i,j)；
2. 以该对为直径端点作圆，并逐个检查所有顶点到圆心的距离，
   给出顶点是否全部落在圆盘内，考虑数值容差 tau；
3. 若最远点对不唯一，则对每组直径圆分别判断是否覆盖。
请先给思路，再给可运行 Python 代码。不要用暴力 O(k^2) 枚举。
\end{lstlisting}

\paragraph{AI 输出摘要}
\begin{lstlisting}[style=AI]
# 旋转卡壳：对凸多边形求直径（省略边界情形处理）
def polygon_diameter(V):
    n = len(V); if n < 3: return 0, (0,0)
    i0 = 0; j = 0; best = -1; pair = (0,0)
    while i0 < 2*n:  # 对踵点收缩
        ...
    return best, pair

# 直径圆覆盖检验：圆心 C、半径 r
def check_cover(V, C, r, eps=1e-9):
    return all(dist(v, C) <= r + eps for v in V)
\end{lstlisting}

\paragraph{人工修改与核验}：AI 给出的思路框架正确，但初始对踵点的选法在
退化多边形下会出现索引越界；团队在循环前增加了“去除重复顶点、去掉自交”
的处理。随后用暴力 $O(k^2)$ 枚举对 6 顶点算例交叉比对，旋转卡壳结果
与暴力结果一致，区域直径均为 $40.8745\,\mathrm{m}$；并针对跨越 $0^\circ$
角的示向、平行边界、无界交集、退化输入分别编写单元测试，均通过。最终
实现的输出与论文式 \texttt{eq:p1-cover} 一致。

\subsection{示例二：问题四多点无信号排除证书的几何论证}

\paragraph{用户提问}
\begin{quote}
在问题四，定向源只有在发射方向前半平面且距离不超过接收半径时才能被
检测到。现某一频道在若干位置 $p$ 获得“无信号”反馈。当存在个数不少于三、
且在 \(1\,\mathrm{m}\) 两两距离的无信号点集合时，能否构造一个排除区域，
使该区域内不存在可能的定向源？请先论证：对任意发射方向，三点凸包
内是否存在一个点不能被三点中的任何一点同时看见，再给出集合运算实现。
\end{quote}

\paragraph{AI 输出}：利用三点凸包与三个以 $1000\,\mathrm{m}$ 为半径的圆盘的
交，即论文式中的 $\mathcal{D}_{ijk}$。对任意过该区域内点 $x$ 的闭半平面，
必至少包含三点中某一点，因此 $x$ 至少与该一点距离不超过
$1000\,\mathrm{m}$，从而与“该点处有源且当前方向无法看见”的假设矛盾。

\paragraph{人工核验}：该几何推理与团队用画图、数值验证的结论一致。由于
关乎“不能凭一次无信号排除源”的关键保证，团队又用穷举采样在
$\mathcal{D}_{ijk}$ 内随机取 $10^5$ 点、枚举 $72$ 个朝向做数值反证，
确认无反例后才采纳；并在实现中对排除区域向内收缩一容差 $\delta$，
避免数值边界误判。此项由人工最终校核。

\subsection{示例三：问题三滚动时域开放路径优化}

\paragraph{用户提问}
\begin{quote}
在每个停驻时刻，我们从当前位置出发，用最近邻加二边交换模拟退火求一条
开放路径，对候选任务点集
candidates。请实现：最近邻初始解、固定随机种子 2-opt 交换，并只返回
并执行整条序列中的第一条行动。注意不能一次性执行完整路径，因为需要
在每步获得新观测后更新规划。
\end{quote}

\paragraph{AI 输出}：给出一个 2-opt 交换加指数退火实现，首行动返回正确。

\paragraph{人工修改}：AI 初版对“终点固定”与“开放路径”的边界处理不当，
团队将其修正为真正的开放端；并加入“每条路径只执行首行动、观测后立即
更新”的控制逻辑，这一滚动规划是团队的核心设计。AI 主要降低了退火初解
的编写工作量，路径正确性依赖团队自己的收敛性测试。

%--------------------- 4. 采纳与核验 ------------------------
\section{对 AI 输出的采纳、人工修改与核验情况（语言润色除外）}

本节按“采纳、修改后采纳、未采纳”三类集中介绍。

\subsection{实现层面}
\begin{itemize}
  \item \textbf{采纳}：可直接复用的函数封装，包括旋转卡壳求直径、
        2-opt 交换、matplotlib 基础绘图、部分 LaTeX 公式排版与校验。
  \item \textbf{修改后采纳}：大多数由 AI 生初版、团队人工修补後使用，
        如 0/360 跨角处理、$d_-$、$d_+$ 公式的投影范围收敛、
        $\varepsilon'=1.0051^\circ$ 的舍入符号，整体人工修改约占百分之
        四十。
  \item \textbf{未采纳}：AI 个别建议，如主张用单点无信号直接排除定向
        源可行域，或每次都重新规划一次全局路径，这些与题设约束或核心
        建模假设冲突，故删除未用。
\end{itemize}

\subsection{4.5 数值与验证}
AI 生成的数值结果经人工独立复核：
\begin{itemize}
  \item 问题一：用暴力 $O(k^2)$ 交叉核对旋转卡壳直径；6 个顶点逐一代入
       直径圆半径检验，与论文表一致。
  \item 问题二：候选点 $(525,60)$、$(392.69,82.76)$ 等坐标由人工按候选点
        公式独立复算，误差不超过 $10^{-3}\,\mathrm{m}$。
  \item 问题三、四：最终清除比例、总虚拟时间等关键结论由团队用独立
        采样脚本复核，AI 仅作为过程可视化与辅助可视化，不负责后者
        确性。
\end{itemize}

\subsection{文献与规范}
参考文献由团队依据题述与规范人工整理，AI 生成的引用线索仅作查询入口，
最终以正式规范为准，规避不合理引用的影响。

\end{document}